% Дополнительное программное исследование: новые вещественные модели.
% Оно не меняет целочисленные сертификаты и аппаратные контрольные точки.
\subsection{Дополнительная проверка обобщения на новых парах адресов}
\label{sec:pair-stage2}
Для проверки чувствительности результатов к разбиению и сравнения с ещё одной
аддитивной моделью обучены новые модели на фиксированном разбиении 211\,043 строк
исходного файла TON\_IoT. Оба направления обмена между неупорядоченной парой
IP-адресов источника и назначения относятся к одной группе; номера портов в
ключ группы не входят. Фиксированное правило SHA-256 с солью распределяет группы
по обучающей, валидационной и тестовой частям в пропорции 60:20:20. Получены
119\,334, 53\,622 и 38\,087 строк соответственно. Пересечений пар адресов и полных
строк из 44 полей между частями нет. Разбиение не обеспечивает независимости
хостов, времени или входных признаков: обучающая и тестовая части имеют 50 общих
хостов, а преобразованные входы 2\,889 тестовых строк (7{,}585\%) встречаются в
обучающей части. В тесте 3\,669 нормальных и 34\,418 атакующих строк; примеры
backdoor, ransomware и XSS отсутствуют.

Все модели используют десять числовых признаков, выбранных по бинарной метке только обучающей части,
и четыре категориальных поля; предобработка обучалась только на обучающей части.
Сравнивались аддитивная KAN степени восемь, дерево глубины пять, MLP с 16
скрытыми нейронами и логистическая GAM с аддитивными кубическими B-сплайнами,
пятью узлами на числовой признак и без взаимодействий. Неизвестные категории
кодировались нулевым one-hot блоком; необученные элементы категории ноль в KAN
обнулены без изменения обучающих логитов. Перед серией выполнен пилот с одним
seed только на валидации. Его оценки не использовались для выбора модели,
гиперпараметров, seed или порога. В основной серии использованы seed
20260916--20260920 и фиксированный seed предобработки 20260916. Все двадцать
моделей сохранены и зафиксированы хешами до вычисления тестовых предсказаний;
порог решения равен 0{,}5. Это новые вещественные модели, отличные от
зафиксированных целочисленных экспортов и аппаратных контрольных точек.

\begin{table}[!htbp]
\centering
\caption{Дополнительный тест с непересекающимися парами адресов:
среднее $\pm$ выборочное стандартное отклонение по пяти seed на одних и тех же
38\,087 строках. Отклонение описывает вариацию оптимизации на одном разбиении,
а не доверительный интервал. FPR и TPR приведены в процентах.}
\label{tab:pair-stage2}
\small
\resizebox{\textwidth}{!}{%
\begin{tabular}{lccccc}
\hline
Модель & BA & F1 & FPR (\%) & TPR (\%) & AUROC\\
\hline
Аддитивная KAN & $0.89484\pm0.00022$ & $0.98695\pm0.00009$ & $20.594\pm0.050$ & $99.563\pm0.019$ & $0.98488\pm0.00027$\\
DT5 & $0.68510\pm0$ & $0.96620\pm0$ & $62.687\pm0$ & $99.707\pm0$ & $0.49014\pm0$\\
MLP16 & $0.98885\pm0.00062$ & $0.99798\pm0.00015$ & $2.044\pm0.152$ & $99.815\pm0.039$ & $0.99806\pm0.00044$\\
Сплайновая GAM & $0.97003\pm0$ & $0.99561\pm0$ & $5.724\pm0$ & $99.730\pm0$ & $0.99606\pm0$\\
\hline
\end{tabular}}
\end{table}

На этом разбиении KAN превосходит DT5 по сбалансированной точности, но уступает
MLP16 и сплайновой GAM (табл.~\ref{tab:pair-stage2}). Значение
F1=0{,}98695 сопровождается FPR=20{,}594\%, что ограничивает практическую
пригодность детектора. Для 35\,198 тестовых строк с преобразованными входами,
отсутствующими в обучающей части, средние показатели KAN составили BA=0{,}89638,
F1=0{,}98613 и FPR=20{,}248\%; в этой подгруппе 3\,623 нормальные и
31\,575 атакующих строк. Подгруппа совпадающих входов содержит лишь 46
нормальных строк. Результаты не доказывают преимущества аддитивной структуры
для обнаружения и не измеряют завышение прежних оценок из-за утечки:
изменились разбиение, состав классов, обученная предобработка и сами модели.
Они уточняют область вклада статьи: проверяемый целочисленный экспорт при
заранее заданном требовании аддитивного представления оценки. Пять seed не
являются независимыми датасетами; результаты пилота и нового теста не
объединяются с аппаратной выборкой.
